\section{Additional parameter-landscape analyses}
\label{supp:s5-parameter}

This note provides additional figures supporting the parameter-landscape results in the main text. We include diagnostics that (i) validate symmetry constraints in the theory-validation regime, (ii) disentangle components of the information-density effect, (iii) close a finite-size discrepancy at high information density ($k=500$), and (iv) illustrate how local coupling $\beta$ induces branch-selection bias and $(q,a)$ coupling in the extended model.

\begin{figure}[H]
  \centering
  \includegraphics[width=0.92\linewidth]{s4_symmetric_diagonal.pdf}
  \caption{Symmetry diagnostics in the theory-validation regime. Along the symmetric constraint, the theoretical constructions satisfy the expected neutrality conditions used to derive the analytic critical point and the landscape map.}
  \label{fig:supp:symmetric-diagonal}
\end{figure}

\begin{figure}[H]
  \centering
  \includegraphics[width=0.92\linewidth]{s4_k_effect_split.pdf}
  \caption{Decomposing the information-density effect. We report additional diagnostics separating how $k$ shapes sensitivity $\chi$ and, through $\chi$, the predicted critical point $r_c$.}
  \label{fig:supp:k-effect-split}
\end{figure}

\begin{figure}[H]
  \centering
  \includegraphics[width=0.92\linewidth]{s4_k500_finite_size.pdf}
  \caption{Finite-size closure for $k=500$. Apparent discrepancies at high information density are resolved by local scanning and finite-size extrapolation, supporting the interpretation that the deviation is a finite-$N$ effect rather than a violation of the analytic prediction.}
  \label{fig:supp:k500-finite-size}
\end{figure}

\begin{figure}[H]
  \centering
  \includegraphics[width=0.92\linewidth]{s4_beta_branch_bias.pdf}
  \caption{Branch-selection bias induced by local coupling. When local coupling $\beta>0$ is introduced, branch selection can become biased, yielding an effective transition behavior that differs from the symmetric well-mixed benchmark.}
  \label{fig:supp:beta-branch-bias}
\end{figure}

\begin{figure}[H]
  \centering
  \includegraphics[width=0.92\linewidth]{s4_beta_q_a.pdf}
  \caption{Joint $(q,a)$ effects under local coupling. Local coupling shifts the macroscopic boundary and modifies the relationship between polarization and activity, illustrating how topology and local reinforcement jointly determine proximity to criticality.}
  \label{fig:supp:beta-q-a}
\end{figure}

